\documentclass[a4paper,12pt]{report}

% Packages needed
\usepackage[utf8]{inputenc}
\usepackage[a4paper, margin=2.5cm]{geometry}
\usepackage{graphicx}
\usepackage{float}
\usepackage{booktabs}
\usepackage{listings}
\usepackage{xcolor}
\usepackage{hyperref}
\usepackage{capt-of}
\usepackage{titlesec}
\usepackage{wrapfig}
\usepackage{needspace}

% Reduce the large empty space before an unnumbered chapter title
\titleformat{\chapter}[display]
  {\normalfont\huge\bfseries}{}{0pt}{\Huge}
\titlespacing*{\chapter}{0pt}{0pt}{30pt}

% Code listing style (Python)
\lstset{
  language=Python,
  basicstyle=\ttfamily\small,
  keywordstyle=\color{blue}\bfseries,
  commentstyle=\color{gray}\ttfamily,
  stringstyle=\color{red}\ttfamily,
  numbers=left,
  numberstyle=\tiny\color{gray},
  stepnumber=1,
  numbersep=5pt,
  frame=single,
  breaklines=true,
  breakatwhitespace=false,
  showstringspaces=false,
  tabsize=2,
  captionpos=b
}

\begin{document}

\chapter*{Single Image Classification Implementation}
\addcontentsline{toc}{chapter}{Single Image Classification Implementation}
\label{chap:implementation}

\section{Overview}
\label{sec:overview}

This chapter presents the implementation of the proposed single-image
classification system based on Convolutional Neural Networks (CNNs). The
implementation covers the complete development pipeline, beginning with
dataset preparation and image preprocessing, followed by the design of the
CNN architecture, model training, performance evaluation, and deployment.
The objective is to develop an accurate and lightweight classification model
suitable for integration into the proposed intelligent workflow automation
platform.

\section{Dataset Preparation}
\label{sec:dataset-preparation}

\begin{wrapfigure}[19]{r}{0.38\textwidth}
    \centering
    \vspace{-45pt}
    \includegraphics[width=0.36\textwidth]{images/dataset_organization.png}
    \caption{Dataset organization and augmentation process.}
    \label{fig:dataset-organization}
\end{wrapfigure}

A labeled image dataset containing multiple object categories was used to
train the proposed CNN model. Before training, the dataset was divided into
three subsets: training, validation, and testing. This separation enables
the model to learn representative visual features during training while
providing an unbiased evaluation on previously unseen images.

To improve the model's generalization capability, data augmentation
techniques were applied to the training images. Operations such as image
rotation, horizontal flipping, and zooming generated additional image
variations without increasing the dataset size, helping reduce overfitting
and improve robustness.

\section{Image Preprocessing}
\label{sec:image-preprocessing}

Before being fed into the CNN, every image undergoes a preprocessing
pipeline to ensure consistent input quality. First, all images are resized
to a fixed resolution compatible with the network architecture. Pixel
values are then normalized to improve numerical stability during
optimization. Finally, class labels are converted into numerical indices
that can be processed by the learning algorithm.

Applying the same preprocessing steps during both training and inference
guarantees consistent model behavior and reliable predictions.

\section{CNN Architecture}
\label{sec:cnn-architecture}

The proposed model employs a lightweight Convolutional Neural Network
designed to extract meaningful visual representations while maintaining
computational efficiency. The architecture consists of multiple
convolutional layers responsible for learning hierarchical image features,
followed by max-pooling layers that gradually reduce spatial dimensions and
computational complexity.

\begin{lstlisting}[caption={Proposed CNN architecture}, label={lst:cnn-architecture}]
model = models.Sequential([
    layers.Conv2D(32, (3,3), activation='relu', input_shape=(128,128,3)),
    layers.MaxPooling2D(2,2),

    layers.Conv2D(64, (3,3), activation='relu'),
    layers.MaxPooling2D(2,2),

    layers.Conv2D(128, (3,3), activation='relu'),
    layers.MaxPooling2D(2,2),

    layers.Flatten(),

    layers.Dense(128, activation='relu'),
    layers.Dropout(0.3),

    layers.Dense(train_gen.num_classes, activation='softmax')
])
\end{lstlisting}

After feature extraction, fully connected layers perform high-level
reasoning, while the final Softmax layer produces the probability
distribution over all object classes. The predicted class corresponds to
the category with the highest probability.

This architecture enables the network to progressively learn low-level
features such as edges and textures before identifying more complex
semantic patterns required for accurate image classification.

The layers of the proposed architecture are described below:

\begin{itemize}
    \item \textbf{Input Layer:} Receives the input RGB image and prepares
        it for feature extraction.
    \item \textbf{Convolution 1:} Extracts low-level features such as
        edges, corners, and simple textures.
    \item \textbf{Max Pooling 1:} Reduces the spatial dimensions while
        preserving the most important features.
    \item \textbf{Convolution 2:} Learns more complex visual patterns by
        combining low-level features.
    \item \textbf{Max Pooling 2:} Further reduces the feature map size and
        computational complexity.
    \item \textbf{Convolution 3:} Extracts high-level semantic features
        that represent meaningful object characteristics.
    \item \textbf{Max Pooling 3:} Compresses the feature maps while
        retaining the most informative features.
    \item \textbf{Flatten Layer:} Converts the extracted feature maps into
        a one-dimensional feature vector.
    \item \textbf{Fully Connected Layer:} Learns the relationship between
        extracted features and object classes to perform classification.
    \item \textbf{Dropout Layer:} Randomly deactivates neurons during
        training to reduce overfitting and improve model generalization.
    \item \textbf{Output Layer (Softmax):} Computes the probability of
        each class and predicts the class with the highest probability as
        the final output.
\end{itemize}

\begin{figure}[H]
    \centering
    \vspace{-5pt}
    \includegraphics[width=0.95\textwidth]{images/cnn_architecture.png}
    \vspace{-5pt}
    \caption{Proposed CNN architecture.}
    \label{fig:cnn-architecture}
\end{figure}

\section{Model Training}
\label{sec:model-training}

The CNN model was trained using the Adam optimizer due to its adaptive
learning capability and fast convergence. Sparse Categorical Crossentropy
was selected as the loss function because the classification problem
involves multiple mutually exclusive classes represented by integer
labels.

To reduce overfitting, an Early Stopping strategy continuously monitored
the validation loss and automatically restored the best-performing model
whenever no further improvement was observed. This approach improves
generalization while avoiding unnecessary training iterations.

The complete training pipeline consists of loading mini-batches, performing
forward propagation, computing the classification loss, updating the
network parameters through backpropagation, and repeating this process
until convergence.

\needspace{15\baselineskip}
\section{Model Evaluation}
\label{sec:model-evaluation}

\begin{wrapfigure}[16]{r}{0.4\textwidth}
    \centering
    \vspace{-35pt}
    \includegraphics[width=0.38\textwidth]{images/confusion_matrix.png}
    \caption{Confusion matrix showing class-wise prediction performance on the test set.}
    \label{fig:confusion-matrix}
\end{wrapfigure}

Following training, the proposed CNN was evaluated using an independent
testing dataset to assess its generalization capability. Several
complementary evaluation metrics were employed to obtain a comprehensive
understanding of the classifier's performance.

Accuracy measures the overall percentage of correctly classified images.
Precision evaluates the reliability of positive predictions, while Recall
measures the ability of the model to identify all relevant instances. The
F1-score combines Precision and Recall into a single balanced metric,
particularly useful when class distributions are uneven. A Confusion
Matrix was generated to visualize class-wise prediction performance and
identify commonly confused categories. Together, these evaluation metrics
provide a comprehensive assessment of both the overall classification
accuracy and the performance of individual classes.

\section{Model Deployment}
\label{sec:model-deployment}

After training and evaluation, the final CNN model was prepared for
deployment. The trained network was saved in HDF5 format (\texttt{.h5}),
while the mapping between numerical class indices and class names was
stored separately in a JSON file. This separation simplifies deployment and
enables predictions to be converted into human-readable labels.

During inference, an input image follows the same preprocessing pipeline
used during training before being passed through the trained CNN. The
network computes the probability for each class, and the class with the
highest probability is selected as the final prediction.

This deployment pipeline ensures that the model behaves consistently
during real-world inference while maintaining the accuracy achieved during
training.

\end{document}

